\par\Needspace{4\baselineskip}
\begin{center}
\hypertarget{native-062-article045}{}\pdfbookmark[2]{45.}{article-045}45.
\end{center}
As already noted, the second part of the first table appended to this work concerns hyperbolic motion. For the argument $A$, which is common to both parts of the table, it gives the logarithm of $B$ and the quantity $C$ to seven decimal places, with the leading digits omitted, and the quantity $T$ first to five, then to six figures. Like the first part, this part extends to $A=0{,}300$, which corresponds to $T=0{,}241207$, $u=2{,}930$ or $=0{,}341$, and $F=\pm52^0\,19'$. Extending it further would have been superfluous (Art. 36).

We now give the order of calculation both for determining the time from the true anomaly and for determining the true anomaly from the time. In the first problem, we obtain $T$ by the formula $T=\frac{e-1}{e+1}\operatorname{tang}\tfrac{1}{2}v^2$. From $T$ our table gives $\log B$ and $C$, whence $A=\frac{(1+C)T}{1-\tfrac{4}{5}T}$. Finally, formula 2) of the preceding article gives $t$. In the second problem, first compute the logarithms of the constants
\[
\begin{aligned}
\alpha&=\frac{75k\BareRoot{(\tfrac{1}{5}+\tfrac{9}{5}e)}}{2q^{\frac{3}{2}}}\mbox{\HistoricalDisplayNoteMark}\\[.7em]
\beta&=\frac{5e-5}{1+9e}\\[.7em]
\gamma&=\BareRoot{\frac{5e+5}{1+9e}}\mathord{\raisebox{.35ex}{.}}
\end{aligned}
\]
\noindent\HistoricalDisplayNoteText{Manuscript note concerning a table that \textsc{Gauss} calculated in 1816 for his own use. The table is still unpublished and will be discussed again in this volume:] in our table $\alpha=\frac{k\BareRoot{(\tfrac{1}{5}+\tfrac{9}{5}e)}}{q^{\frac{3}{2}}}\mathord{\raisebox{.35ex}{.}}$}
\SourceLeaf{063}Then determine $A$ from $t$ in exactly the same way as for elliptic motion. The mean motion $\frac{\alpha t}{B}$ corresponds to the true anomaly $w$ in \textsc{Barker}'s table, and we have $A=\beta\operatorname{tang}\tfrac{1}{2}w^2$. First find an approximate value of $A$, either neglecting the factor $B$ or estimating it if the means are available. Our table then supplies an approximate value of $B$ with which to repeat the operation. The new value of $B$ so obtained will hardly ever need a perceptible correction, so another iteration will be unnecessary. Once $A$ has been corrected, take $C$ from the table. This gives
\[
\begin{aligned}
\operatorname{tang}\tfrac{1}{2}v&=\frac{\gamma\operatorname{tang}\tfrac{1}{2}w}{\BareRoot{(1+\tfrac{4}{5}A+C)}},\\[.7em]
r&=\frac{(1+\tfrac{4}{5}A+C)q}{(1-\tfrac{1}{5}A+C)\cos\tfrac{1}{2}v^2}\mathord{\raisebox{.35ex}{.}}
\end{aligned}
\]
It follows that the formulas for elliptic and hyperbolic motion agree exactly if we treat $\beta$, $A$, and $T$ as negative quantities in hyperbolic motion.
